%% Question bank chapter file - generated by bank_tools.py build (do not edit while the bank is still being built:
%% it is regenerated from the per-test data in Question Bank\data). Every question is one bankq block:
%%   \begin{bankq}[id=..., source={...}, paper=.., qno=.., type=..., subtopic={...}, answer={...}, answer_src=key|solved,
%%                 working={...}, difficulty=1-10, diff_parts={K S I C F}, diff_note={...},
%%                 group=<passage id or empty>, also_in={duplicate copies in other tests}]
%% difficulty = min(10, K+S+I+C+F-1): K concepts 1-3, S steps 1-4, I insight 0-2, C computation 0-2, F format 0-1
%%   ...question LaTeX, options, figures (\bfig / \bfigi -> img/)...
%%   \end{bankq}
%% Passages shared by several questions are bankpassage blocks; their questions carry group=<passage id>.
\documentclass[10pt]{article}
\usepackage[a4paper,margin=1.8cm]{geometry}
\usepackage[T1]{fontenc}
\usepackage{amsmath,amssymb,mathtools}
\usepackage{graphicx}
\usepackage[export]{adjustbox}
\usepackage{xcolor,needspace,keyval}
\usepackage{enumitem}
\usepackage{booktabs,array,tabularx}
\usepackage[hidelinks]{hyperref}
\setlength{\parindent}{0pt}
\setlength{\parskip}{3pt}
\makeatletter
\define@key{bq}{id}{\def\bq@id{#1}}\define@key{bq}{source}{\def\bq@src{#1}}\define@key{bq}{paper}{\def\bq@paper{#1}}
\define@key{bq}{qno}{\def\bq@qno{#1}}\define@key{bq}{type}{\def\bq@type{#1}}\define@key{bq}{subtopic}{\def\bq@sub{#1}}
\define@key{bq}{answer}{\def\bq@ans{#1}}\define@key{bq}{answer_src}{\def\bq@asrc{#1}}\define@key{bq}{working}{\def\bq@work{#1}}
\define@key{bq}{group}{\def\bq@group{#1}}\define@key{bq}{also_in}{\def\bq@also{#1}}
\define@key{bq}{difficulty}{\def\bq@diff{#1}}\define@key{bq}{diff_parts}{\def\bq@dparts{#1}}\define@key{bq}{diff_note}{\def\bq@dnote{#1}}
\newenvironment{bankq}[1][]{\def\bq@work{}\def\bq@also{}\def\bq@group{}\def\bq@diff{?}\def\bq@dparts{}\def\bq@dnote{}\setkeys{bq}{#1}%
  \par\needspace{4\baselineskip}\medskip\noindent{\small\textbf{\bq@id}\hfill\color{gray}[\bq@src, Paper \bq@paper, Q\bq@qno]
  \bq@type\ $\cdot$ \bq@sub\ $\cdot$ \textbf{D\bq@diff}}\par\nopagebreak\relax}
  {\par{\small\color{teal}Answer: \textbf{\bq@ans}\ifx\bq@work\@empty\else\ (solved: \bq@work)\fi
  \ifx\bq@also\@empty\else\par\color{gray}Also in: \bq@also\fi}\par}
\define@key{bp}{id}{\def\bp@id{#1}}\define@key{bp}{source}{\def\bp@src{#1}}\define@key{bp}{questions}{\def\bp@qs{#1}}
\newenvironment{bankpassage}[1][]{\setkeys{bp}{#1}\par\needspace{4\baselineskip}\bigskip\noindent
  {\small\textbf{Passage \bp@id}\hfill\color{gray}[\bp@src; questions \bp@qs]}\par\nopagebreak\relax\itshape}{\par}
\makeatother
\newcommand{\bfig}[2]{\par\smallskip{\centering\includegraphics[width=#1pt,max width=0.62\linewidth,max height=0.32\textheight]{img/#2}\par}\smallskip}
\newcommand{\bfigi}[2]{\includegraphics[width=#1pt,max width=\linewidth,max height=0.32\textheight,valign=c]{img/#2}}
\newcommand{\bmissing}[1]{\textit{[figure missing in the source conversion: \detokenize{#1}]}}
\newenvironment{qoptions}{\begin{itemize}[leftmargin=1.8em,itemsep=1pt,topsep=2pt,label={}]}{\end{itemize}}
\newcommand{\opt}[2]{\item[(#1)] #2}
\providecommand{\fwstart}{}\providecommand{\fwend}{}\providecommand{\qfullsplit}{}
\newcolumntype{L}[1]{>{\raggedright\arraybackslash}p{#1}}
\newcolumntype{C}[1]{>{\centering\arraybackslash}p{#1}}
\newcolumntype{R}[1]{>{\raggedleft\arraybackslash}p{#1}}
\providecommand{\sectionhdr}[1]{}\providecommand{\subsectionhdr}[1]{}


\begin{document}
{\Large\bfseries Metallurgy --- Question Bank}\par\smallskip
{\small Chemistry chapter C32. Questions from all converted test papers, grouped by source series. Answers: printed key where it matches the paper, otherwise solved (answer\_src=solved, with working). Exact repeats across papers appear once, with also\_in listing the other copies.}\par\medskip\hrule

\section*{Score-I papers}
\begin{bankq}[id=SI12-P1-CHEM-S1-Q08, source={Score-I - Test 12 (09 Mar 2022)}, paper=1, qno=8, type=multi, subtopic={Refining}, answer={B,C}, answer_src=key, difficulty=3, diff_parts={K1 S2 I0 C0 F1}, diff_note={Aqueous electrorefining fails for metals more reactive than hydrogen (Na, Al)}]
In electrofining of metals anode and cathode are taken as thick slab of impure metal and strip of pure metal respectively while the electrolyte is aqueous solution of a complex metal salt. This method cannot be applied for the refining of :
\begin{qoptions}
\opt{A}{Copper}
\opt{B}{Sodium}
\opt{C}{Aluminium}
\opt{D}{Silver}
\end{qoptions}
\end{bankq}

\section*{TOAS papers}
\begin{bankq}[id=TOAS06-P1-CHEM-S2-Q11, source={TOAS - Test 06 (17 Apr 2022)}, paper=1, qno=11, type=multi, subtopic={Ores}, answer={A,C,D}, answer_src=key, difficulty=3, diff_parts={K1 S2 I0 C0 F1}, diff_note={Identify iron-free minerals: rutile, cerussite, pyrolusite}]
Which of the following minerals do/does not contain iron ?
\begin{qoptions}
\opt{A}{Rutile}
\opt{B}{Ilmenite}
\opt{C}{Cerrusite}
\opt{D}{Pyrolusite}
\end{qoptions}
\end{bankq}
\begin{bankq}[id=TOAS06-P1-CHEM-S3-Q06, source={TOAS - Test 06 (17 Apr 2022)}, paper=1, qno=6, type=numerical, subtopic={Extraction}, answer={3}, answer_src=solved, working={Soluble in excess NaOH(aq): hydrated alumina (amphoteric), SiO$_2$ and ZnO (amphoteric); Fe$_2$O$_3$, TiO$_2$, CuO, MgO, NiO are not: 3}, difficulty=4, diff_parts={K1 S3 I0 C0 F1}, diff_note={Which oxides dissolve in excess NaOH: amphoteric and acidic ones}]
A sample of ore contains\\
(a) iron oxide\\
(b) hydrated alumina\\
(c) $\mathrm{SiO_2}$\\
(d) $\mathrm{TiO_2}$\\
(e) CuO\\
(f) MgO\\
(g) NiO\\
(h) ZnO\\
Number of components soluble in excess $\mathrm{NaOH_{aq.}}$
\end{bankq}
\begin{bankq}[id=TOAS08-P1-CHEM-S3-Q01, source={TOAS - Test 08 (15 May 2022)}, paper=1, qno=1, type=numerical, subtopic={Extraction}, answer={6.00}, answer_src=key, difficulty=5, diff_parts={K1 S4 I0 C0 F1}, diff_note={Match eight processes and ores with the metals concerned}]
Find number of process or ore given in Column-I are correctly represented with Column-II. You have to consider P for P, Q for Q, R for R etc.
\begin{center}\small\renewcommand{\arraystretch}{1.1}\setlength{\extrarowheight}{3.5pt}\setlength{\tabcolsep}{3pt}
\begin{tabular}{|c|l|c|l|}
\hline
\multicolumn{2}{|c|}{\parbox[t]{3.4cm}{\centering (I)\\ Process/Ore}} & \multicolumn{2}{c|}{\parbox[t]{3.4cm}{\centering (II)\\ Related to}}\\[3pt]
\hline
(P) & Matte smelting & (P) & Cu\\[3pt]
(Q) & Cyanide leaching & (Q) & Au\\[3pt]
(R) & Mond's process & (R) & \parbox[t]{2.3cm}{Reaction of metal $+\mathrm{H_2}$}\\[3pt]
(S) & Fused oxide electrolysis & (S) & Al\\[3pt]
(T) & Chromite ore & (T) & Fe $+$ Cr\\[3pt]
(U) & Van Arkel process & (U) & Zr\\[3pt]
(V) & Bacterial leaching & (V) & Low grade Cu ore\\[3pt]
(W) & Separation of $\mathrm{FeWO_4}$ & (W) & Extraction of Pb\\[3pt]
\hline
\end{tabular}
\end{center}
\end{bankq}
\begin{bankq}[id=TOAS08-P1-CHEM-S3-Q02, source={TOAS - Test 08 (15 May 2022)}, paper=1, qno=2, type=numerical, subtopic={Extraction}, answer={2.00}, answer_src=key, difficulty=6, diff_parts={K1 S4 I1 C0 F1}, diff_note={Seven statements on metallurgy to be checked}]
Consider the following statements\newline
(1) $\mathrm{CO_2}$ (g) $+$ C (s) $\to$ 2CO(g) ; it is an exothermic reaction\newline
(2) Ilmenite is an carbonate ore containing two metals\newline
(3) In froth floatation process KCN act as an activator when galena contains impurity of ZnS\newline
(4) Liquation is applicable for the purification of Sn , Pb, Zn\newline
(5) Au $+$ NaCN $+$ $\mathrm{O_2}$ $+$ $\mathrm{H_2O}$ $\to$ $\mathrm{Na[Au(CN)_4]}$ reaction occurs during leaching of gold\newline
(6) $\mathrm{Al_2O_3}$ can be obtained in solution from red bauxite by using concentrated NaOH\newline
(7) MgO was reduced by C at very high temperature followed by a rapid cooling with the help of inert gas\newline
The number of correct statement(s) in the above given statements are :
\end{bankq}
\begin{bankq}[id=TOAS08-P2-CHEM-S1-Q02, source={TOAS - Test 08 (15 May 2022)}, paper=2, qno=2, type=multi, subtopic={Extraction}, answer={B,C,D}, answer_src=key, difficulty=5, diff_parts={K1 S4 I0 C0 F1}, diff_note={Four statements on iron extraction}]
Which of the following statements(s) about extraction of Fe is/are correct?
\begin{qoptions}
\opt{A}{Cast iron is made by mixing pig iron with scrap iron and coke using air blast at 300K}
\opt{B}{By oxidizing impurities of cast iron in a furnace lined with Haematite, wrought iron can be produced}
\opt{C}{Layer of Molten slag prevents further oxidation state of Molten metal layer in Blast furnace}
\opt{D}{Pig iron has more carbon content than that of cast iron}
\end{qoptions}
\end{bankq}
\begin{bankq}[id=TOAS08-P2-CHEM-S3-Q01, source={TOAS - Test 08 (15 May 2022)}, paper=2, qno=1, type=numerical, subtopic={Extraction}, answer={5.00}, answer_src=key, difficulty=5, diff_parts={K1 S4 I0 C0 F1}, diff_note={Reactions of extraction with gaseous products}]
Number of reactions in which atleast one gaseous product is formed.\newline
(i) $\mathrm{Al_2O_3+SiO_2+N_2+C}\xrightarrow{1800^\circ\mathrm{C}-2000^\circ\mathrm{C}}$\newline
(ii) $\mathrm{4Au+8NaCN+O_2+2H_2O}\rightleftharpoons$\newline
(iii) $\mathrm{SnSiO_3+CaO+C}\to$\newline
(iv) $\mathrm{2Cu_2S+3O_2}\to$\newline
(v) $\mathrm{XeF_2+H_2O}\to$\newline
(vi) $\mathrm{ZnO+C}$ (coke) $\xrightarrow{1673\mathrm K}$
\end{bankq}
\begin{bankq}[id=TOAS09-P1-CHEM-S2-Q07, source={TOAS - Test 09 (30 May 2022)}, paper=1, qno=7, type=multi, subtopic={Ores}, answer={B,D}, answer_src=key, difficulty=6, diff_parts={K2 S3 I1 C0 F1}, diff_note={Ores: standard approach for this subtopic; 2 concept(s), 3 step(s), needs a non-obvious insight}]
Which of the following ore(s) contain(s) copper metal?
\begin{qoptions}
\opt{A}{Dolomite}
\opt{B}{Chalcopyrite}
\opt{C}{Fool's gold}
\opt{D}{Malachite}
\end{qoptions}
\end{bankq}
\begin{bankq}[id=TOAS09-P2-CHEM-S1-Q02, source={TOAS - Test 09 (30 May 2022)}, paper=2, qno=2, type=single, subtopic={Ores}, answer={C}, answer_src=key, difficulty=2, diff_parts={K1 S2 I0 C0 F0}, diff_note={Ores: standard approach for this subtopic; 1 concept(s), 2 step(s)}]
Calamine malachite, magnetite and cryolite, respectively are:
\begin{qoptions}
\opt{A}{ZnSO$_4$, CuCO$_3$, Fe$_2$O$_3$, AlF$_3$}
\opt{B}{ZnSO$_4$, Cu(OH)$_2$, Fe$_3$O$_4$, Na$_3$AlF$_6$}
\opt{C}{ZnCO$_3$, CuCO$_3\cdot$Cu(OH)$_2$, Fe$_3$O$_4$, Na$_3$AlF$_6$}
\opt{D}{ZnCO$_3$, CuCO$_3$, Fe$_2$O$_3$, Na$_3$AlF$_6$}
\end{qoptions}
\end{bankq}
\end{document}
